\section{Existence statements, and the type of a number}\label{ch15:sec:exists}
To prove an existence statement, we name a witness and check that it has the required property (Section~\ref{ch04:sec:directions}). To prove that some natural number $n$ satisfies $n+1=4$, we offer the witness $3$ and check that $3+1=4$. In Lean:
\par\noindent\begin{minipage}{\linewidth}
\begin{lstlisting}
theorem a_witness : Exists (fun n : Nat => n + 1 = 4) := by
  exact Exists.intro 3 rfl
\end{lstlisting}
\end{minipage}\par

Lean writes $\exists n\in\NN:\ n+1=4$ as \code{Exists (fun n : Nat => n + 1 = 4)}, which the infoview displays as $\exists$\,\code{n, n + 1 = 4}. The input to \code{Exists} is a function that turns each natural number $n$ (that is what \code{n : Nat} says) into the proposition $n+1=4$, true for $n=3$ and false for $n=2$. The whole expression says that this proposition is true for at least one natural number.

To prove the statement, \code{Exists.intro} takes two inputs: a witness, here $3$, and a proof that the witness has the property, which is a proof of \code{3 + 1 = 4}. That proof is \code{rfl}, used this time inside an expression rather than as a tactic; there it names the proof that an object equals itself. Lean accepts it because \code{3 + 1} reduces to \code{4}. With the witness $2$, Lean would ask for a proof of the false equation \code{2 + 1 = 4}, and \code{rfl} would fail.

To use a hypothesis that something exists, we must first name the object, so that we can refer to it, and record its property, as a paper proof does when it says ``let $n$ be such a number, so that $n+1=4$''; Section~\ref{ch16:sec:surj} does this with the tactic \code{cases}.

\subsection*{The type of a number decides what subtraction means}
The witness $3$ had to be a natural number, because the statement declared \code{n : Nat}. The type also decides what the arithmetic symbols mean. A natural number cannot be negative, so Lean's natural-number subtraction is \emph{truncated} at zero: \code{a - b} is the usual difference if $a\ge b$, and $0$ if $a<b$. The notation \code{(e : T)} tells Lean to read \code{e} as an object of type \code{T}, and so decides which subtraction is meant: \code{(7 - 10 : Nat)} is $0$, while \code{(7 - 10 : Int)} is $-3$, because \code{Int}, Lean's name for $\ZZ$, contains negative numbers. To see this, put the line \code{\#eval (7 - 10 : Nat)} in your file, outside any proof: the command \code{\#eval} asks Lean to compute an expression, and the result \code{0} appears in the infoview when the cursor is on that line. The line \code{\#eval (7 - 10 : Int)} gives \code{-3}.

Truncated subtraction breaks familiar laws: $(a-b)+b=a$ holds for all integers, but in \code{Nat} it fails whenever $a<b$, since \code{(7 - 10) + 10} is $0+10=10$. So an argument that subtracts a larger number from a smaller one cannot be copied unchanged into natural-number subtraction. This is one reason why the statement contract of Habit~2 (Section~\ref{sec:magnitude}) records the domain of every variable. When an assistant offers to repair a formal proof, check whether it has quietly changed the kind of number being used.
